\documentclass[conference]{IEEEtran}

% Codificación y configuración general
\usepackage[utf8]{inputenc}
\usepackage[T1]{fontenc}
\usepackage[spanish,es-nodecimaldot]{babel}

% Paquetes habituales para artículos IEEE
\usepackage{cite}
\usepackage{amsmath,amssymb,amsfonts}
\usepackage{graphicx}
\usepackage{textcomp}
\usepackage{xcolor}
\usepackage{url}
\usepackage{booktabs}
\usepackage{listings}
\usepackage{algorithm}
\usepackage{algorithmic}

% TikZ se conserva para que pueda habilitarse cuando sea necesario.
\usepackage{tikz}
\usetikzlibrary{shapes.geometric,positioning}

\def\UrlBreaks{\do\/\do-\do\_}

\lstset{
    basicstyle=\ttfamily\small,
    showstringspaces=false,
    breaklines=true,
    frame=single,
    numbers=left,
    numberstyle=\tiny\color{gray},
    inputencoding=utf8,
    extendedchars=true,
    literate={á}{{\'a}}1 {é}{{\'e}}1 {í}{{\'i}}1 {ó}{{\'o}}1
    {ú}{{\'u}}1 {ñ}{{\~n}}1 {ü}{{\"u}}1 {Á}{{\'A}}1 {É}{{\'E}}1
    {Í}{{\'I}}1 {Ó}{{\'O}}1 {Ú}{{\'U}}1 {Ñ}{{\~N}}1 {Ü}{{\"U}}1
}

% Descomenta y ajusta esta ruta si el documento utiliza imágenes.
% \graphicspath{{../img/}}

\title{Título del trabajo}

\author{
    \IEEEauthorblockN{Melchor Alexander Araiza Zavala}
\IEEEauthorblockA{
\texttt{ma.araizazavala@ugto.mx}
}
\and
\IEEEauthorblockN{Leonardo Daniel Aviña Neri}
\IEEEauthorblockA{
% Licenciatura en Ingeniería de Datos e Inteligencia Artificial\\
% Universidad de Guanajuato\\
\texttt{ld.avinaneri@ugto.mx}
}
\and
\IEEEauthorblockN{Ricardo Ignacio Pérez Mendoza}
\IEEEauthorblockA{
\texttt{ri.perezmendoza@ugto.mx}
}
\and
\IEEEauthorblockN{Regina Madahí Rodríguez Jasso}
\IEEEauthorblockA{
\texttt{rm.rodriguezjasso@ugto.mx}
}
\and
\IEEEauthorblockN{Emmanuel Vargas Cisneros}
\IEEEauthorblockA{
\texttt{e.vargascisneros@ugto.mx}
}
}

\begin{document}
\maketitle

\begin{abstract}
Lorem ipsum dolor sit amet, consectetur adipiscing elit. Integer tincidunt,
nisl at facilisis consequat, justo sem porta neque, vitae tincidunt massa
libero non ipsum. Este resumen debe presentar brevemente el problema, la
metodología, los resultados principales y la conclusión del trabajo.
\end{abstract}

\begin{IEEEkeywords}
palabra clave 1, palabra clave 2, metodología, resultados, análisis
\end{IEEEkeywords}

\section{Introducción}
Lorem ipsum dolor sit amet, consectetur adipiscing elit. Quisque consequat,
urna sit amet fermentum ullamcorper, justo mauris posuere risus, vitae
convallis lectus erat vel arcu. En esta sección se presenta el contexto del
problema, su importancia y la motivación del trabajo.

El objetivo general del estudio es describir una propuesta reproducible para
analizar el problema seleccionado. Los objetivos específicos pueden incluir
la identificación de variables, el diseño de una solución y la evaluación de
los resultados obtenidos.

\section{Marco teórico}
Lorem ipsum dolor sit amet, consectetur adipiscing elit. Suspendisse potenti.
Praesent eget lectus sed sem commodo consequat. Las definiciones, conceptos
fundamentales y trabajos relacionados deben colocarse en esta sección.

Una expresión matemática puede escribirse de la siguiente manera:
\begin{equation}
    f(x) = \sum_{i=1}^{n} w_i x_i,
\end{equation}
donde $w_i$ representa un peso y $x_i$ representa una variable de entrada.

\section{Metodología}
Lorem ipsum dolor sit amet, consectetur adipiscing elit. Donec vitae
sollicitudin mauris. Nam euismod, massa quis consectetur tincidunt, sem
libero tincidunt lectus, at commodo justo neque quis erat.

\subsection{Materiales y herramientas}
Indique aquí los datos, herramientas, lenguajes de programación, bibliotecas
y recursos empleados en el desarrollo.

\subsection{Procedimiento}
Describa aquí las etapas necesarias para reproducir el trabajo. Puede
organizar el procedimiento mediante una lista, un algoritmo o subsecciones.

\begin{algorithm}[htbp]
\caption{Procedimiento de ejemplo}
\label{alg:procedimiento}
\begin{algorithmic}[1]
\STATE Recibir los datos de entrada.
\STATE Procesar la información de acuerdo con el método propuesto.
\STATE Evaluar y registrar los resultados.
\RETURN Resultados del análisis.
\end{algorithmic}
\end{algorithm}

\section{Resultados}
Lorem ipsum dolor sit amet, consectetur adipiscing elit. Curabitur ac
sollicitudin justo. Los resultados deben presentarse de forma clara,
acompañados de tablas, figuras o métricas cuando sea necesario.

\begin{table}[htbp]
\centering
\caption{Tabla de resultados de ejemplo.}
\label{tab:resultados}
\begin{tabular}{lcc}
\toprule
Categoría & Métrica 1 & Métrica 2 \\
\midrule
Caso A & 0.00 & 0.00 \\
Caso B & 0.00 & 0.00 \\
\bottomrule
\end{tabular}
\end{table}

% Gráfico de ejemplo con TikZ. Descomenta este bloque cuando lo necesites.
% \begin{figure}[htbp]
% \centering
% \begin{tikzpicture}[
%     block/.style={
%         rectangle,
%         rounded corners=3pt,
%         draw=black,
%         align=center,
%         minimum width=0.75\columnwidth,
%         minimum height=0.7cm
%     },
%     arrow/.style={-stealth, draw=black},
%     node distance=0.4cm
% ]
% \node[block] (entrada) {Datos de entrada};
% \node[block, below=of entrada] (proceso) {Procesamiento};
% \node[block, below=of proceso] (salida) {Resultados};
% \draw[arrow] (entrada) -- (proceso);
% \draw[arrow] (proceso) -- (salida);
% \end{tikzpicture}
% \caption{Diagrama de flujo de ejemplo.}
% \label{fig:diagrama}
% \end{figure}

\section{Discusión}
En el mundo natural, las variables de temperatura y humedad no operan como entidades aisladas o independientes, sino que existe una interacción termodinámica directa entre ambas. A medida que la temperatura de un espacio se incrementa, la capacidad del aire para retener vapor de agua cambia dinámicamente. Esto provoca que, ante un ascenso de temperatura constante, la percepción del calor se intensifique al mismo tiempo que la humedad en el ambiente tiende a decrecer de la misma manera.
Este fenómeno se manifiesta en situaciones extremas donde la humedad no puede sostenerse en estado gaseoso en el aire caliente, o bien requiere de una superficie con menor temperatura o un punto donde se mantenga fresco o con una fuente de agua cercana, para que pueda condensarse y manifestarse en el entorno. De la misma forma, en escenarios de baja temperatura o alta saturación de humedad, la evaporación se reduce, modificando la percepción térmica global. La representación gráfica simultánea de ambas variables de entrada permite auditar si los límites definidos para cada etiqueta lingüística capturan adecuadamente este comportamiento de la realidad.
Al proyectar valores específicos dentro de las funciones de membresía de ambas variables, la discusión técnica puede ir más allá de la simple evaluación de una fórmula individual, permitiendo examinar cómo interactúa el espacio de entradas en conjunto. Por ejemplo, al evaluar casos donde la temperatura se desplaza hacia la etiqueta CALIENTE, el modelo difuso debe estar estructurado para interpretar si esa condición coincide con una reducción en la pertenencia hacia el conjunto de humedad ALTA y un ajuste hacia humedad MEDIA o BAJA. Del mismo modo, el uso de funciones trapezoidales en los extremos acota los estados de saturación y límites físicos del entorno, evitando lecturas fuera de rango y dejando claro que la temperatura no puede llegar a ser indefinidamente más calurosa ni tan fría. La posibilidad de visualizar el grado de activación de cada etiqueta en un instante dado garantiza que el sistema no esté considerando combinaciones físicas contradictorias o, dicho de otra forma, que lo que presente se acerque a la vida real. 
Asimismo, la presencia de regiones de traslape en las funciones trapezoidales y triangulares de ambas variables lingüísticas (como en la transición entre TEMPLADA y CALIENTE) asegura que las variaciones no sean procesadas mediante saltos discontinuos. En la naturaleza, el paso de un ambiente fresco y húmedo a uno cálido y seco ocurre mediante un cambio continuo y no abrupto (evitando cambios irreales de calor a frío). La evaluación visual de entradas de prueba confirma que la lógica emula dicho comportamiento: conforme el valor de temperatura avanza sobre la pendiente de su respectiva curva, la transición gradual en los grados de pertenencia refleja el comportamiento con el que cambian las condiciones atmosféricas en la realidad. 
En conclusión, la integración de esta representación gráfica dentro del análisis no solo valida la corrección matemática al fuzzificar, sino que permite verificar que la arquitectura difusa responde a las leyes físicas del ambiente, asegurando que las reglas que se activen posteriormente en el motor de inferencia no disparen antecedentes incompatibles y tomen decisiones coherentes con el comportamiento termodinámico real.


\section{Conclusiones}
En el desarrollo de este trabajo se logró implementar un sistema difuso basado en la matriz BIOFAM y en la regla de selección por valor de activación máximo, demostrando que los resultados obtenidos se aproximan a las condiciones térmicas reales del problema. Para el caso de la temperatura y la humedad, el comportamiento alcanzado respondió a las expectativas teóricas, logrando resolver la problemática inicial dentro del tiempo planteado. En cuanto a la experiencia durante el desarrollo del ejercicio, el proceso permitió comprender de manera práctica la transición entre la modelación matemática de las funciones de membresía y la toma de decisiones lógica ante variables con incertidumbre.
Esta vivencia técnica puso en evidencia la importancia de un adecuado diseño en las etiquetas lingüísticas, optando por mantener únicamente tres conjuntos por variable (fresco/frío, templado/medio, cálido/alto) con el fin de ofrecer una representación simple, pero a su vez clara, para evitar la ambigüedad que surge de la percepción térmica individual al aumentar o decrecer la temperatura y los procesos de aclimatación del cuerpo humano según la región geográfica. La principal limitante del sistema radica en la dependencia de una correcta introducción de datos por parte del usuario, quiere decir que al colocar los valores, estos sean correctos y no superen la escala planteada. No obstante, este factor es habitual en aplicaciones de ingeniería, lo que reafirma que estos sistemas deben ser operados por personal capacitado y que todo modelo generalizado requiere de una etapa continua de pruebas experimentales para validar y refinar sus resultados.


\section*{Agradecimientos}
Se acredita un agradecimiento al docente de la materia Sánchez Yáñez por su valiosa orientación y por facilitar un código de ejemplo, que sirvió como punto de partida para el desarrollo de esta investigación. Su aportación inicial no solo sentó los cimientos técnicos necesarios para plantear formalmente el problema, sino que también proporcionó la claridad metodológica indispensable para abordar el proyecto. En particular, la guía brindada para la formulación e interpretación de la matriz BIOFAM, así como la pauta para la estructuración del sistema de reglas, resultando en elementos para adquirir un entendimiento preciso, del problema que se buscaba resolver.
De igual manera, otorgamos un merecido reconocimiento a cada uno de los integrantes que conformaron este equipo de trabajo. El éxito de este proyecto es el reflejo directo del compromiso, la dedicación y el esfuerzo colaborativo mantenido a lo largo de todas las etapas del desarrollo. Cada aportación individual representó una contribución sustancial al logro del objetivo final: desde las labores iniciales de análisis, abstracción y delimitación del problema, hasta el proceso de definición, ajuste y validación de cada una de las reglas. Todos y cada uno de estos pasos, resultaron indispensables para consolidar una solución.


\begin{thebibliography}{00}
\bibitem{referencia1}
Autor, ``Título de la referencia,'' \textit{Nombre de la revista o
conferencia}, vol. 1, no. 1, pp. 1--10, 2026.

\bibitem{referencia2}
Autor, \textit{Título del libro}, Editorial, 2026.
\end{thebibliography}

\end{document}
